\documentclass{article}
\usepackage{amsmath, amssymb, amsthm}
\usepackage{hyperref}
\usepackage{geometry}
\geometry{margin=1in}

\title{Erd\H{o}s Problem \#1184}
\date{Last updated: 2026-04-04}
\author{Source: erdosproblems.com}

\begin{document}
\maketitle

\noindent\textbf{Status:} OPEN \quad \textbf{No prize}

\noindent\textbf{Tags:} number theory, primes

\medskip

\noindent\textbf{Problem Statement:}

Let $f(n,k)$ count the number of $1\leq i\leq k$ such that $P(n+i)&gt;k$ (where $P(m)$ is the largest prime divisor of $m$). Is it true that, if $\alpha&gt;1$ is such that $n=k^{\alpha+o(1)}$, then\[f(n,k)=(1-\rho(\alpha)+o(1))k,\]where $\rho$ is the Dickman function ?




Erd\H{o}s \cite{Er76e} proved that, for every $\alpha&#62;1$, if $k$ is sufficiently large and $n&#62;k^\alpha-k$ then\[f(n,k) &#62; \left(1-\frac{1}{\alpha}+c_\alpha\right)k\]for some constant $c_\alpha&#62;0$, and if $1&lt;\alpha&lt;2$ and $n\leq k^\alpha-k$ then\[f(n,k) &lt; (\alpha-1+o(1))k.\]He knew of no non-trivial bounds when $\alpha \geq 2$. Ramachandra, Shorey, and Tijdeman \cite{RST75b} proved that if $n&gt;\exp(c(\log k)^2)$ for a constant $c&#62;0$ then $f(n,k)\geq k-\pi(k)$.




References

[Er76e] Erd\H{o}s, P., Problems and results on consecutive integers . Publ. Math. Debrecen (1976), 271-282.

[RST75b] Ramachandra, K. and Shorey, T. N. and Tijdeman, R., On {G}rimm's problem relating to factorisation of a block of consecutive integers. {II} . J. Reine Angew. Math. (1976), 192--201.

\noindent\textbf{Related OEIS sequences:} possible


\bigskip
\noindent\small{Source: \url{https://www.erdosproblems.com/1184}}
\end{document}
